\documentclass[10pt,a4paper]{amsart}
\usepackage[margin=19mm]{geometry}
\usepackage{amsmath,amssymb,mathtools}
\usepackage[scr=rsfs,cal=euler]{mathalfa}
\usepackage{xfrac}
\usepackage[unicode,hidelinks,bookmarks=false]{hyperref}
\newcommand{\ie}{i.e.\ }
\newcommand{\cf}{cf.\ }
\newcommand{\resp}{respectively\ }
\newcommand{\Subsection}{\subsection}
\providecommand{\nobreakdash}{\nobreak\hbox{-}}
\setlength{\emergencystretch}{2em}
\multlinegap0pt
\usepackage{amssymb}
\usepackage{caption}
\usepackage[shortlabels]{enumitem}
\usepackage{tikz}
\usepackage[all]{xy}
\usepackage{xcolor}
\newtheorem{lemme}{Lemme}
\newtheorem{proposition}{Proposition}
\DeclareMathOperator{\Frob}{Frob}
\DeclareMathOperator{\Hom}{Hom}
\newcommand\Qp{\mathbb{Q}_{p}}
\DeclareMathOperator{\Sel}{Sel}
\newcommand\Zp{{\mathbb{Z}_{p}}}
\newcommand\bF{\mathbb{F}}
\newcommand\bQ{\mathbb{Q}}
\newcommand\cE{\mathcal{E}}
\newcommand\cO{\mathcal{O}}
\newcommand{\ob}[1]{\mkern 1.5mu\overline{\mkern-1.5mu#1\mkern-1.5mu}\mkern 1.5mu}
\title{Théorie d'Iwasawa des motifs d'Artin et des formes modulaires de poids 1}
\author{Alexandre Maksoud}
\date{Source: arXiv:1811.05368v3 (CC BY 4.0)}
\begin{document}
\maketitle

\noindent\emph{Source and licence.} Théorie d'Iwasawa des motifs d'Artin et des formes modulaires de poids 1. Version arXiv:1811.05368v3 (CC BY 4.0), distributed by arXiv under the Creative Commons Attribution 4.0 International licence, http://creativecommons.org/licenses/by/4.0/, as recorded in the arXiv licence metadata for this exact version and shown on the abstract page https://arxiv.org/abs/1811.05368v3. Full licence notice: \texttt{licenses/arxiv-1811.05368-CC-BY-4.0.txt}.\par\smallskip

\noindent\textit{Editorial scope.} Counted unit: the article's proposition prop:cyc (source0.tex lines 1187--1189). The article's complete original proof of it is delivered with it (source0.tex lines 1190--1236, one \texttt{proof} environment of 47 lines). No mathematics is written, repaired or re-derived: the statement and the proof are the article's own lines, in the article's own notation, and no retained line is substituted or re-flowed.  Retained uncounted as the source-local context the proof needs: lines 923--928; lines 930--946; lines 950--958; lines 960--962.  Nothing was omitted from any retained span, so the package carries no omission record.  No retained line contains a citation, so the wrapper carries no bibliography and no bibliography entry.  No retained line contains a figure environment, an image-inclusion command or drawing code, so no image is delivered and the package declares no asset dependency.  Editorial formatting only, in addition: an AMS \texttt{amsart} preamble that restates the article's own declarations and macros used by the retained lines, namely the environments \texttt{lemme}, \texttt{proposition} on separate counters, together with the standard packages those lines need.  The retained environments print in this document as Lemme 1, Proposition 1 in order of appearance, whereas the article prints the same items with the numbering of its own full document.  The complete native source of the article as distributed in the arXiv e-print for this exact version is bundled unchanged as \texttt{sources/101960/source0.tex}.
\begin{equation}\label{isomQbar}
\left\{ \begin{array}{rcl}
\ob{\bQ}_p[G_n] & \longrightarrow & \Hom_\textrm{ct}(U_n,\ob{\bQ}_p) \\
g          & \mapsto         & \left( x\otimes c \mapsto \log_p\left(g^{-1}(x)\otimes c \right) \right).
\end{array} \right.
\end{equation}

\begin{lemme}\label{lem:description_morphismes_Qpbar}
	Soit $f \in \Hom_{G_n}(U_n,V_{p,\chi})$. Alors $f$ s'écrit sous la forme 
	$$f(x\otimes c) = \sum_{g \in G_n} \log_p \left(g^{-1}(x)\otimes c \right)\left(\rho\otimes\chi\right)(g)(\vec{v}) = \lambda_p(x\otimes c) \cdot \vec{v}, $$ 
	pour un unique vecteur $\vec{v}\in V_{p,\chi}$, et la restriction à $U_{n,v}$ est donnée par 
	$$f(u) = \sum_{g\in G_{n,v}}  \log_p\left(g^{-1}(u)\right)\left(\rho\otimes\chi\right)(g)(\vec{v}) = \lambda_p(u)\cdot \vec{v}. $$
	De plus, $f(U_{n,v}) \subseteq V_{p,\chi}^+$ si et seulement si $\vec{v}\in V_{p,\chi}^+$.
\end{lemme}
\begin{proof}
	On a $\Hom_{G_n}(U_n,V_{p,\chi}) = \left(\Hom_{\textrm{ct}}(U_n,\ob{\bQ}_p) \otimes V_{p,\chi} \right)^{G_n}$, donc d'après l'isomorphisme (\ref{isomQbar}) il existe des vecteurs $\vec{v}_g \in V_{p,\chi}$ tels que 
	$$f(x\otimes c) = \sum_{g \in G_n} \log_p \left(g^{-1}(x)\otimes c \right)\vec{v}_g.$$
	En utilisant la $G_n$-invariance, on a facilement $\vec{v}_g = \left(\rho\otimes\chi\right)(g)(\vec{v_1})$. 
	On a donc bien $f(x\otimes c) = \lambda_p(x\otimes c) \cdot \vec{v}$ avec $\vec{v}=\vec{v_1}$. 
	L'unicité de $\vec{v}$ est claire, car l'image de $\lambda_p$ engendre $\ob{\bQ}_p[G_n]$. 
	La description de $f$ sur $U_{n,v}$ se déduit de celle de $\lambda_p$. 
	Enfin, comme $\lambda_p(\ob{\bQ}_p U_{n,v})=\ob{\bQ}_p[G_{n,v}]$, on a aussi $f(\ob{\bQ}_p U_{n,v}) = \ob{\bQ}_p[G_{n,v}] \cdot \vec{v}$. 
	Or, $V_{p,\chi}^+$ est $G_{n,v}$-stable, donc $f(U_{n,v}) \subseteq V_{p,\chi}^+$ si et seulement si $\vec{v}\in V_{p,\chi}^+$.
\end{proof}

\begin{lemme}\label{lem:calcul_matriciel}
	Fixons $1\leq i\leq d^+$. Pour tout $g\in G_n$, on a une égalité de matrices lignes à $d$ entrées : 
	$$\begin{pmatrix} \log_p\left(g(\epsilon_{i,1})\right) & \cdots &\log_p\left(g(\epsilon_{i,d})\right)\end{pmatrix}=\begin{pmatrix}
	s_{i,1} & \cdots &  s_{i,d}
	\end{pmatrix} \left(\rho\otimes\chi\right)(g).$$
\end{lemme}
\begin{proof}
	Soit $g\in G_n$. Notons $a_{j,k}$ les coefficients de $\rho\otimes\chi(g)$ dans la base $(\vec{v}_1,\ldots,\vec{v}_d)$. Pour tout $1\leq j\leq d$, on a $g(\epsilon_{i,j})=\epsilon_{i,1}^{a_{1,j}} \ldots \epsilon_{i,d}^{a_{d,j}}$. En prenant le logarithme $p$-adique et en concaténant les égalités pour $1\leq j\leq d$, on retrouve la formule souhaitée.
\end{proof}

\begin{proposition}\label{prop:S^+}
	L'application $\beta_{n,\chi}$ est injective si et seulement si la matrice carrée $S^+:=(s_{i,j})_{1\leq i,j\leq d^+} \in M_{d^+}(\ob{\bQ}_p)$ est inversible. 
\end{proposition}

\begin{proposition}\label{prop:cyc}
	Le $\cO$-module $\Sel^\#$ est isomorphe au quotient $T_p^+/p^{-1}S^+\left(T_p^+\right)$.
\end{proposition}

\begin{proof}
	Le choix de la base de $T_p$ (adaptée à $T_p^+$ et diagonalisant $\Frob_v$) identifie $T_p^-$ avec un supplémentaire de $T_p^+$ dans $T_p$, \textit{i.e.}, $T_p=T_p^+\bigoplus T_p^-$. On a de même $V_p=V_p^+ \bigoplus V_p^-$, ainsi que $D_p=D_p^+ \bigoplus D_p^-$. 
	On note $(f^+,f^-)$ les coordonnées d'un morphisme $f$ à valeurs dans $V_p$ relativement à cette décomposition. Si $f\in \Hom_G(U,V_p)$, alors le Lemme \ref{lem:description_morphismes_Qpbar} permet d'écrire $f$ sous la forme 
	$$f(x\otimes c) = \sum_{g \in G} \log_p \left(g^{-1}(x)\otimes c \right)\rho(g)(\vec{v}) = \lambda_p(x\otimes c) \cdot \vec{v}, $$ 
	pour un unique vecteur $\vec{v}\in V_p$, et pour tout $x\otimes c \in U \subseteq \left(\cO_H \otimes \Zp\right)^\times$. La restriction à $U_{v}$ est donnée par 
	\begin{equation}\label{eq:expression_f_locale}
	f(u) = \sum_{g\in G_{v}}  \log_p\left(g^{-1}(u)\right)\rho(g)(\vec{v}) = \lambda_p(u)\cdot \vec{v}. 
	\end{equation}
	
	\begin{lemme}
		\'Ecrivons $\vec{v}=\vec{v}^+ \oplus \vec{v}^-$ selon la décomposition $V_p=V_p^+\bigoplus V_p^-$. On a $f^\pm(U_v) \subseteq T_p^\pm$ si et seulement si $p \vec{v}^\pm \in T_p^\pm$. En particulier, on a $f(U)\subseteq T_p$ si et seulement si $p\vec{v}\in T_p$.
	\end{lemme}
	\begin{proof}[Preuve du Lemme]
		 Pour $1\leq i \leq d$, notons $f_i:U_v \longrightarrow L$ la $i$-ème coordonnée de $f$, notons $v_i \in L$ celle de $\vec{v}$ et notons $\zeta_i\in L$ la $i$-ème valeur propre de $\Frob_v$ dans la base ambiante $\{\vec{t}_i\}_i$. Il nous suffit de montrer que, pour tout $1\leq i \leq d$, on a $f_i(U_v)\subseteq \cO$ si et seulement si $pv_i\in \cO$. Fixons $1\leq i \leq d$. L'expression (\ref{eq:expression_f_locale}) pour $f$ montre que l'on a $f_i(u)=\sigma_i\circ \log_p(u) v_i$ pour $u\in U_v=1+p\cO_{H_v}$, où $\sigma_i$ est l'application $\Zp$-linéaire $$\sigma_i :  \left\{\begin{array}{ccl} H_v & \longrightarrow & L \\
		  x&\mapsto & \sum_{j=0}^{\#G_v-1}\Frob^{j}_v(x)\zeta_i^{-j}. 
		  \end{array}\right.$$ 
		 Notons que $\sigma_i$ est effectivement à valeurs dans $L$ car on a supposé que $H_v\subseteq L$. Comme $H_v/\Qp$ est non-ramifiée, le logarithme $p$-adique envoie bijectivement $U_v$ sur $p\cO_{H_v}$ et donc $f_i(U_v)=pv_i\sigma_i\left(\cO_{H_v}\right)$. Pour montrer l'équivalence, il suffit donc de montrer que $\sigma_i\left(\cO_{H_v}\right)\subseteq \cO$ mais que $\sigma_i\left(\cO_{H_v}\right) \not\subseteq p\cO$. Autrement dit, il suffit de voir que $\sigma_i\left(\cO_{H_v}\right)\subseteq \cO$ et qu'il existe $x_0\in\cO_{H_v}$ tel que $\sigma_i(x_0)\in\cO^\times$. La première affirmation est claire car $\zeta_i$ est une racine de l'unité (puisque $\Frob_v\in G$ est d'ordre fini) et donc $\zeta_i\in \cO$. Pour la deuxième, on réduit $\sigma_i$ modulo $p$ (qui est une uniformisante de $H_v$ et de $L$) pour obtenir 
		 $$\ob{\sigma}_i :  \left\{\begin{array}{ccl} \bF_{H_v} & \longrightarrow & \bF_{L} \\
		 \bar{x}&\mapsto & \sum_{j=0}^{\#G_v-1}\zeta_i^{-j}\bar{x}^{p^j}, 
		 \end{array}\right.$$ 
		 où $\bF_{H_v}$ et $\bF_{L}$ sont les corps résiduels respectifs de $H_n$ et de $L$. On doit vérifier que $\ob{\sigma}_i$ n'est pas identiquement nulle. Or, ceci est vrai pour toute application polynômiale sur $\bF_{H_v}$ de degré strictement inférieur à $\#\left(\bF_{H_v}\right)=p^{\#G_v}$, et donc c'est vrai en particulier pour $\ob{\sigma}_i$.
	\end{proof}
	On termine la preuve de la Proposition \ref{prop:cyc}. Soit $f$ comme précédemment. D'après le Lemme \ref{lem:calcul_matriciel} et la Proposition \ref{prop:S^+}, et avec les notations de la preuve de la Proposition \ref{prop:S^+}, on a $f(\cE_p)\subseteq T_p$ si et seulement si les matrices $M_i$ de taille $d\times d$ sont à coefficients dans $\cO$ pour tout $1\leq i \leq d^+$. Or, $\rho$ étant absolument irréductible, $M_i$ est la matrice d'une homothétie de rapport $\lambda_i=\frac{\#G}{d} \left(s_{i,1}v_1+\ldots + s_{i,d}v_d\right)$, où $(v_1,\cdots,v_d)$ sont les coordonnées de $\vec{v}$ dans la base fixée de $T$. Comme $p\nmid \#G \cdot d$ d'après \textbf{(deg)} et \textbf{(dim)}, et comme $s_{i,j}\in p\cO$, on a les équivalences : 
	
	$$\begin{array}{rcl}
	\left\{\begin{array}{l}
	f(U_{v}) \subseteq V_p^+ +T_p \\ f(\cE_p)\subseteq T_p 
	\end{array}\right. &\Longleftrightarrow & \left\{\begin{array}{l}
	f^-(U_{v}) \subseteq T_p^- \\ \forall 1\leq i \leq d^+, \  M_i\in M_d(\cO) 
	\end{array}\right. \\
	&\Longleftrightarrow & \left\{\begin{array}{l}
	p\vec{v}^-\in T_p^- \\ \forall 1 \leq i \leq d^+, \ s_{i,1}v_1+\ldots + s_{i,d}v_{d}\in \cO
	\end{array}\right. \\
	&\Longleftrightarrow & \left\{\begin{array}{l}
	pv_{d^++1},\cdots,pv_d\in \cO \\ \forall 1 \leq i \leq d^+, \ s_{i,1}v_1+\ldots + s_{i,d^+}v_{d^+}\in \cO 
	\end{array}\right. \\
	&\Longleftrightarrow & \left\{\begin{array}{l}
	pv_{d^++1},\cdots,pv_d\in \cO \\ S^+ \cdot\vec{v}^+ \in T_p^+ ,
	\end{array}\right. \\
	&\Longleftrightarrow & \left\{\begin{array}{l}
	\vec{v}^-\in p^{-1}T_p^- \\ \vec{v}^+ \in \left(S^+\right)^{-1}(T_p^+).
	\end{array}\right. 
	\end{array}
	$$
	Ainsi, le $\cO$-module $\Sel^\#$ s'identifie à $\left(\left(S^+\right)^{-1}(T_p^+) \times p^{-1}T_p^-\right) /\left(p^{-1}T_p^+ \times p^{-1}T_p^-\right) \simeq T_p^+/p^{-1}S^+\left(T_p^+\right)$.
	
\end{proof} 

\end{document}
